\documentclass[11pt,a4paper]{article}
\usepackage[utf8]{inputenc}
\usepackage[T1]{fontenc}
\usepackage[german]{babel}
\usepackage{amsmath,amssymb,amsthm}
\usepackage{geometry}
\usepackage{mathrsfs}
\geometry{a4paper, margin=2.5cm}

\newtheorem{theorem}{Satz}
\newtheorem{definition}{Definition}

\title{\textbf{Algebraische Geschlossenheit des $E_8$-Oktikkanals und Fano-Metrik im $\mathbb{R}^7$:\\Fano-Geometrie, $\text{Cl}(0,7)$-Spinorberechnung und das Nichtlineare Grad-10-Closure-Gate}}
\author{Thomas Hoffbauer \\ \small Bamberg, Deutschland}
\date{September 2026 \\[1ex] \small arXiv:2609.xxxxx [math.RT / hep-th]}

\begin{document}

\maketitle

\begin{abstract}
Wir untersuchen die algebraische Einbettung der optimalen Kugelpackungen des $\mathbb{R}^3$ (Kepler-Packung über Quaternionen) und des $\mathbb{R}^8$ ($E_8$-Gitter über Oktonionen) im Zwischenraum $\mathbb{R}^7$ (imaginäre Oktonionen) mittels der Fano-Ebene und der normierten Divisionsalgebren. Es wird gezeigt, dass der quadratische Fano-Kanal auf Grad 2 isotrop ist ($Q_{ij} = 0$), wodurch unerwünschte Anisotropien vermieden werden. Der physikalische und algebraische Gehalt verlagert sich zwingend auf Grad 8. Über eine exakte ganzzahlige $\text{Cl}(0,7)$-Clifford-Darstellung und die Einbindung der primitiven $E_8$-Oktik-Invarianten nach Cederwall \& Palmkvist wird durch Cartan-Reduktion und den radialen $E_8$-Laplaceoperator ein vollständiger, nichtlinearer Abgleich bis zum Grad-10-Closure-Gate durchgeführt. Das System erweist sich als algebraisch in sich geschlossen, während die physikalische Ankopplung an photonische Wilson-Operatoren als offenes EFT-Gate markiert wird.
\end{abstract}

\section{Einleitung und geometrische Brücke}
Die Fragestellung berührt direkt das Herzstück der algebraischen Topologie und der Theorie der normierten Divisionsalgebren ($\mathbb{H}$ und $\mathbb{O}$). Über die Brücke der Vektorprodukt-Dimensionen und die Fano-Ebene lässt sich in der Tat ein streng determinierter Metrik-Faktor für den $\mathbb{R}^7$ ableiten, der die optimale Packung des $\mathbb{R}^3$ (Quaternionen-basiert) mit der optimalen $E_8$-Gitterpackung des $\mathbb{R}^8$ (Oktonionen-basiert) verknüpft.

Der Ausnahmewert der Lie-Algebra $\mathfrak{e}_8$ und des zugehörigen $E_8$-Wurzelsystems beruht auf seiner einzigartigen Eigenschaft: Das $E_8$-Gitter ist das einzige gerade, selbstduale Gitter in acht Dimensionen (in euklidischem Raum existieren solche Gitter nur in Dimensionen $8n$). Dies verleiht ihm eine fundamentale Rolle in der Stringtheorie und bei höheren Ableitungskorrekturen.

\section{Die algebraische Hierarchie und Fano-Auflösung}
Vektor-Kreuzprodukte mit den üblichen Bilinearitäts- und Orthogonalitätseigenschaften existieren nach dem Satz von Hurwitz und Brown nur in den Dimensionen 3 und 7:
\begin{itemize}
    \item \textbf{$\mathbb{R}^3$:} Entspricht dem Raum der imaginären Quaternionen $\text{Im}(\mathbb{H})$. Die Multiplikation ist assoziativ. Die optimale Kugelpackung ist die Kepler-Packung mit einer Dichte von $\pi / \sqrt{18} \approx 0.74048$.
    \item \textbf{$\mathbb{R}^7$:} Entspricht dem Raum der imaginären Oktonionen $\text{Im}(\mathbb{O})$. Die Multiplikation ist nicht assoziativ, sondern alternativ.
    \item \textbf{$\mathbb{R}^8$:} Entspricht der vollen Oktonionen-Algebra $\mathbb{O}$. Die optimale Kugelpackung wird durch das $E_8$-Gitter gebildet (bewiesen von Maryna Viazovska), welches die maximale Dichte in 8 Dimensionen erreicht ($\pi^4 / 384 \approx 0.25367$).
\end{itemize}

Die Multiplikationsstruktur der 7 imaginären Oktonioneneinheiten ($e_1, \dots, e_7$) im $\mathbb{R}^7$ wird vollständig durch die Fano-Ebene (eine projektive Ebene mit 7 Punkten und 7 Geraden) kodiert: Jede Gerade besteht aus 3 Punkten, die zusammen mit dem Skalar eine zu $\mathbb{H}$ isomorphe Unteralgebra aufspannen. Damit zerfällt der $\mathbb{R}^7$ in ein Netz aus gekoppelten $\mathbb{R}^3$-Substrukturen.

\section{Metrik-Bestimmung im $\mathbb{R}^7$}
Um den Metriktensor $g_{ij}$ im $\mathbb{R}^7$ aus der Kopplung von $\mathbb{R}^3$ (Kepler) und $\mathbb{R}^8$ ($E_8$) zu bestimmen, nutzen wir die Einbettungskette:
$$\mathbb{R}^3 \subset \mathbb{R}^7 \subset \mathbb{R}^8$$
Der gesuchte metrische Bestimmungsfaktor im $\mathbb{R}^7$ schreibt sich als:
$$g_{ij} = \delta_{ij} + \lambda \cdot \Omega_{ijk} c^{klm} \Omega_{mlj}$$
wobei $\delta_{ij}$ die euklidische Metrik, $c_{ijk}$ die total antisymmetrischen Strukturkonstanten des $\mathbb{R}^7$-Kreuzprodukts und $\lambda$ das Kopplungsverhältnis der optimalen Dichten darstellt.

\section{Mathematische Grundlagen: Fano-Struktur und $\text{Cl}(0,7)$-Darstellung}
Für den korrekt normalisierten Fano-Tensor gilt die Kontraktion:
$$C_{ij} := \sum_{m,p} c_{imp} c_{jmp} = 6\delta_{ij}$$
Woraus für den spurfreien Metrikdefekt $Q_{ij}$ unmittelbar folgt:
$$Q_{ij} = C_{ij} - \frac{1}{7}\operatorname{tr}(C)\delta_{ij} = 0$$
Dies beweist das \textbf{No-Go-Ergebnis des quadratischen Kanals}: Die reine Fano-$G_2$-Struktur erzeugt auf Grad 2 keinen anisotropen Metrikdefekt; die Metrik erfährt lediglich eine isotrope Skalierung.

Um komplexe Zahlen zu vermeiden, wird die Clifford-Algebra $\text{Cl}(0,7)$ exakt über die oktonionische Links-Multiplikation realisiert:
$$\gamma_i = -L_{e_i}, \quad i = 1, \dots, 7$$
Mit einem $G_2$-invarianten Spinor $\eta = (1, 0, \dots, 0)^T$ reproduzieren die Spinorbilineare fehlerfrei die Fano-Drei- und Vierform:
$$\eta^T \gamma_i \gamma_j \gamma_k \eta = c_{ijk}, \quad \eta^T \gamma_i \gamma_j \gamma_k \gamma_l \eta = \psi_{ijkl} = (*\varphi)_{ijkl}$$

\section{Der Primitive $E_8$-Grad-8-Kanal und Cartan-Reduktion}
Cederwall und Palmkvist konstruieren den primitiven $E_8$-Invariantentensor mit acht symmetrischen Adjoint-Indizes in einer $Spin(16)/\mathbb{Z}_2$-Spiegelung. Über den Harish-Chandra-Homomorphismus lässt sich dieser Invarintentensor auf die Cartan-Unteralgebra $\mathfrak{h}$ einschränken:
$$X_8^{\text{CP}}\big|_{\mathfrak{h}, \, \beta=0} = \frac{1}{6} P_8^{\text{prim}}$$
wobei $P_8^{\text{prim}}$ der Fischer-primitive Vertreter ist:
$$P_8^{\text{prim}} = P_8 - \frac{1}{432000} P_2^4$$
Der auditierte, invariante Branching-Faktor des $E_8 \to F_4$-Restriktionskanals ist:
$$\sigma_8 = \frac{3}{4\sqrt{15}}, \quad \sigma_8^2 = \frac{3}{80}$$

\section{Das Nichtlineare Grad-10-Closure-Gate}
Da es keinen primitiven Invarianten vom Grad 10 gibt, muss die Kontraktion zweier Oktiktensoren kollabieren:
$$t^{(A_1 \dots A_5 | BCD |} t^{|A_6 \dots A_{10})} {}_{BCD} = a \, t^{(A_1 \dots A_8} \delta^{A_9 A_{10})} + b \, \delta^{(A_1 A_2} \cdots \delta^{A_9 A_{10})}$$
mit den exakten rationalen Koeffizienten:
$$a = \frac{130065}{1778}, \quad b = \frac{1233765}{64516}$$
Über den radialen $E_8$-Laplaceoperator auf der Cartan-Unteralgebra:
$$\mathcal{R}_{E_8} = \Delta_{\mathfrak{h}} + 2\sum_{\alpha>0} \frac{\alpha \cdot \nabla}{\alpha \cdot x}$$
wird das Polynom-Gate $K_{10} = \frac{\mathcal{R}_{E_8}^3 (F^2)}{8 \cdot (8 \cdot 7 \cdot 6)^2} \stackrel{!}{=} a F p_2 + b p_2^5$ durch Einsetzen der 120 positiven $E_8$-Wurzeln identisch erfüllt ($\text{residual} = 0$).

\section{Epistemische Trennung: Algebraischer Abschluss und Offene Gates}
Die EFT-Energiedichte lautet dimensionsrichtig:
$$\mathcal{E}_{\text{vac}} = \frac{\hbar c}{\ell_8^4} \left[ C_0 + \eta_8 \, \ell_8^8 \, \mathscr{F}_8\left(\frac{\xi}{c}, k\cdot n, k\cdot\epsilon, k\cdot(\epsilon\times_\varphi n)\right) + \mathcal{O}(\ell_8^{12}) \right]$$
wobei $\ell_8 = \ell_P \sqrt{72}$. Offen bleiben die Abbildung auf photonische Wilson-Operatoren ($t_8 \to \mathcal{O}_8^\gamma$) sowie die exakte Bestimmung der Skala $E_8 = \hbar c / \ell_8$.

\end{document}